\section{Related Work}

\paragraph{Agent safety benchmarks.}
ToolEmu evaluates LM agents using an LM-emulated sandbox and highlights risks such as privacy leakage and financial loss in tool-rich settings \cite{toolemu}. AgentDojo provides dynamic environments for prompt-injection attacks and defenses over realistic agent tasks \cite{agentdojo}. These benchmarks motivate tool-use safety as a first-class problem. Our work asks a narrower pre-commit question: before a task succeeds or fails end-to-end, has the candidate action been bound to the effect, resource, operation, authorization, and provenance facts that determine whether it may commit?

\paragraph{Pre-execution and step-level guardrails.}
ToolSafe/TS-Guard studies proactive step-level tool invocation safety and feedback-driven guardrails \cite{toolsafe}. Safiron studies foundational pre-execution guardrails trained with synthetic data and evaluated on Pre-Exec Bench \cite{safiron}. These systems ask whether a planned or candidate action is risky before execution. Our custom stress tests ask a more specific semantic question: whether the monitor's decision is bound to realized effect, resource, authorization, operation mode, and provenance rather than surface form.

\paragraph{Structural prompt-injection defenses and provenance.}
CaMeL separates trusted control from untrusted data and uses capability ideas to prevent unauthorized data flows \cite{camel}. IPIGuard uses a tool dependency graph to reduce unintended tool invocations under indirect prompt injection \cite{ipiguard}. ARGUS constructs provenance-style influence structure for LLM-agent defenses against context-aware prompt injection \cite{weng2026argus}. AgentVisor uses semantic virtualization, privilege separation, and tool-call interception to mediate agent execution \cite{ying2026agentvisor}. These works shift defenses from prompt-only filtering toward execution structure. Tool-effect binding is compatible with that shift, but it targets a different layer: the candidate action must first be decomposed into the atoms whose control source and provenance can be checked.

\paragraph{Capability, authorization, and least privilege.}
Capability systems were introduced as a way to name both authority and object references in computation \cite{dennis1966programming}. The confused-deputy problem shows why ambient authority can let an otherwise legitimate program exercise the wrong authority on behalf of the wrong principal \cite{hardy1988confused}. Progent introduces programmable privilege control policies for LLM agents \cite{progent}. MiniScope reconstructs permission hierarchies and applies a least-privilege framework for authorizing tool-calling agents \cite{miniscope}. CaMeL also uses capability-style structure to separate data and control \cite{camel}. These works are closest to the authorization motivation. Their policies are only as accurate as the object being checked: least privilege over a tool call can still miss an unauthorized Bcc, public-link change, or commit-mode shift if those atoms are not represented.

\paragraph{Mediation before enforcement.}
Classical security design separates the policy to be enforced from the object over which mediation occurs. Saltzer and Schroeder's complete-mediation principle requires that every access to every object be checked against authority \cite{saltzer1975protection}. Information-flow and provenance foundations similarly emphasize that security decisions depend on how information or influence reaches an action, not only on the final command string \cite{denning1976lattice,buneman2001why,cheney2009provenance}. In LLM-agent tool use, the mediated object is often implicit in a tool name and argument string. This paper contributes the semantic binding layer before enforcement: a side-effectful tool call must first be decomposed into effect-resource-operation atoms so that guardrails, capability restrictions, least-privilege policies, or provenance checks can be applied to the correct object.

\paragraph{Positioning.}
Prior work provides benchmarks, pre-execution guards, capability restrictions, least-privilege policies, and provenance-aware structure. Our contribution is the semantic binding layer before enforcement: a side-effectful tool call must first be decomposed into effect-resource-operation atoms so that policy can be applied to the correct object.
